%---
% title: "Đạo hàm và ý nghĩa hình học"
% date: 2026-08-17
% tags: [toan, latex]
% summary: "Tìm hiểu đạo hàm qua góc nhìn hình học: định nghĩa, quy tắc tính và ứng dụng."
%---
\section{Giới thiệu}

Đạo hàm là một trong những khái niệm nền tảng nhất của giải tích.
Nó mô tả tốc độ thay đổi tức thời của một hàm số.

\section{Định nghĩa}

Đạo hàm của hàm số \(f\) tại điểm \(x\) được định nghĩa bởi giới hạn:

\begin{equation}
f'(x) = \lim_{h \to 0} \frac{f(x+h)-f(x)}{h}
\end{equation}

\section{Ví dụ bằng mã}

Dưới đây là một hàm Python tính đạo hàm số:

\begin{lstlisting}[language=Python]
def dao_ham(f, x, h=1e-6):
    return (f(x + h) - f(x - h)) / (2 * h)
\end{lstlisting}

\section{Kết luận}

Đạo hàm giúp ta hiểu độ dốc của đồ thị tại từng điểm, và là công cụ
trung tâm trong tối ưu hóa, vật lý và học máy.
